\begin{titlepage}
\centering
{\Large Notas didácticas de una entrevista a fondo\par}
\vspace{1.0cm}
{\huge\bfseries Diez años de investigación multimodal, la inspiración de o1 y el próximo momento GPT-4}\par
\vspace{0.5cm}
{\Large Zhang Xiaojun Podcast: Li Guangmi (李广密), co-host, conversa con Zhang Xiangyu (张祥雨)\par}
\vspace{0.35cm}
{\large Elaboradas a partir del video público, la estructura de capítulos, la transcripción, la descripción del video y diagramas conceptuales propios\par}
\vspace{0.3cm}
{\large 2025-06-03\par}
\vspace{0.85cm}
\includegraphics[width=0.88\textwidth,height=0.42\textheight,keepaspectratio]{cover.jpg}\par
\vfill
\begin{tcolorbox}[width=0.92\textwidth, colback=black!2!white, colframe=black!55, sharp corners]
\textbf{Título del video}: 102. Conversación con Zhang Xiangyu: las dificultades de la investigación multimodal y los dos «momentos GPT-4» de los próximos dos años\par
\textbf{Canal del video}: Zhang Xiaojun Podcast\par
\textbf{Duración del video}: 02:28:44\par
\textbf{Enlace del video}: \href{\videourl}{\nolinkurl{\videourl}}\par
\textbf{Nota sobre los materiales}: YouTube no ofrece subtítulos descargables; el texto principal se elaboró a partir de una transcripción por bloques con faster-whisper large-v3 (CPU, int8), la descripción del video, la estructura de capítulos y figuras didácticas propias. Las tomas fijas de la entrevista no se repiten en el texto principal.\par
\textbf{Página del proyecto}: \href{\repourl}{\nolinkurl{\repourl}}\par
\end{tcolorbox}
\end{titlepage}

\begingroup
\small
\setlength{\parskip}{0pt}
\renewcommand{\baselinestretch}{0.92}\selectfont
\tableofcontents
\endgroup
\newpage

